% Part: first-order-logic
% Chapter: tableaux
% Section: soundness.tex

\documentclass[../../../include/open-logic-section]{subfiles}

\begin{document}

\iftag{FOL}
      {\olfileid{fol}{tab}{sou}}
      {\olfileid{pl}{tab}{sou}}
      
\olsection{Correctheid}

\begin{explain}
!!^a{derivation}ssysteem, zoals de tableaumethode, is \emph{correct}
als het geen zaken kan !!{derive} die in werkelijkheid niet gelden.
Correctheid is dus een soort gegarandeerde veiligheidseigenschap van
!!{derivation}ssystemen. Afhankelijk van de bewijstheoretische
eigenschap die aan de orde is, willen we bijvoorbeeld weten dat
\begin{enumerate}
\item elke !!{derivable}~$!A$ \iftag{FOL}{geldig}{een tautologie} is;
\item als !!a{sentence} uit enkele andere !!{derivable} is, hij ook een
  gevolg daarvan is;
\item als een verzameling !!{sentence}s inconsistent is, zij
  onvervulbaar is.
\end{enumerate}
Dit zijn belangrijke eigenschappen van !!a{derivation}ssysteem. Als
een ervan niet geldt, is het !!{derivation}ssysteem gebrekkig: het zou
te veel !!{derive}. Het is daarom van het grootste belang de
correctheid van !!a{derivation}ssysteem vast te stellen.

Omdat al deze bewijstheoretische eigenschappen met behulp van een of
ander gesloten !!{tableau} zijn gedefinieerd, moeten we voor het
bewijs van (1)--(3) hierboven iets bewijzen over de semantische
eigenschappen van gesloten !!{tableau}s. We definiëren eerst wat het
betekent dat !!a{signed formula} in een structuur wordt vervuld. Daarna
tonen we aan dat geen enkele structuur alle aannamen van een gesloten
!!{tableau} vervult. (1)--(3) volgen vervolgens als gevolgen van dit
resultaat.
\end{explain}

\begin{defn}
\iftag{FOL}{!!^a{structure}}{!!^a{valuation}}~$\iftag{FOL}{\Struct{M}}{\pAssign{v}}$
\emph{vervult} !!a{signed formula} $\sFmla{\True}{!A}$ dan en slechts
dan als $\iftag{FOL}{\Sat{M}}{\pSat{v}}{!A}$, en vervult
$\sFmla{\False}{!A}$ dan en slechts dan als
$\iftag{FOL}{\Sat/{M}}{\pSat/{v}}{!A}$. $\iftag{FOL}{\Struct{M}}{\pAssign{v}}$
vervult een verzameling !!{signed formula}s~$\Gamma$ dan en slechts
dan als elke $\sFmla{S}{!A} \in \Gamma$ wordt vervuld. $\Gamma$~is
\emph{vervulbaar} als er een \iftag{FOL}{!!a{structure}}{!!a{valuation}}
is die haar vervult, en anders \emph{onvervulbaar}.
\end{defn}

\begin{thm}[Correctheid]
  \ollabel{thm:tableau-soundness}
  Als $\Gamma$ een gesloten !!{tableau} heeft, is $\Gamma$
  onvervulbaar.
\end{thm}

\begin{proof}
We noemen een tak van !!a{tableau} vervulbaar dan en slechts dan als
de verzameling !!{signed formula}s erop vervulbaar is. We noemen
!!a{tableau} vervulbaar als het ten minste één vervulbare tak bevat.

We tonen het volgende aan: als een vervulbaar !!{tableau} met een van
de afleidingsregels wordt uitgebreid, levert dit altijd een
vervulbaar !!{tableau} op. Daarmee is de stelling bewezen: elk gesloten
!!{tableau} ontstaat door afleidingsregels toe te passen op het
!!{tableau} dat uitsluitend uit aannamen uit~$\Gamma$ bestaat. Als
$\Gamma$ dus vervulbaar was, zou elk !!{tableau} ervoor vervulbaar
zijn. Een gesloten !!{tableau} is echter duidelijk niet vervulbaar:
elke tak bevat zowel $\sFmla{\True}{!A}$ als
$\sFmla{\False}{!A}$, en geen structuur kan~$!A$ tegelijk wel en niet
vervullen.

Stel dat we een vervulbaar !!{tableau} hebben, d.w.z. !!a{tableau} met
ten minste één vervulbare tak. De toepassing van een afleidingsregel
voegt !!{signed formula}s aan een tak toe of splitst een tak in tweeën.
Als het !!{tableau} een vervulbare tak heeft die door de betrokken
regeltoepassing niet wordt uitgebreid, blijft deze in het uitgebreide
!!{tableau} een vervulbare tak. Het uitgebreide tableau is dan dus
vervulbaar. We hoeven daarom alleen het geval te beschouwen waarin een
regel op een vervulbare tak wordt toegepast.

Laat $\Gamma$ de verzameling !!{signed formula}s op die tak zijn en
laat $\sFmla{S}{!A} \in \Gamma$ de !!{signed formula} zijn waarop de
regel wordt toegepast. Als de regel geen gesplitste tak oplevert,
moeten we aantonen dat de uitgebreide tak, d.w.z. $\Gamma$ samen met
de conclusies van de regel, nog steeds vervulbaar is. Als de regel wel
een gesplitste tak oplevert, moeten we aantonen dat ten minste een van
de twee resulterende takken vervulbaar is.
  
Eerst beschouwen we de mogelijke inferenties die geen gesplitste tak
opleveren.
\begin{enumerate}
\item De tak wordt uitgebreid door $\TRule{\True}{\lnot}$ toe te
  passen op $\sFmla{\True}{\lnot !B} \in \Gamma$. De uitgebreide tak
  bevat dan de !!{signed formula}s $\Gamma \cup
  \{\sFmla{\False}{!B}\}$. Stel dat
  $\iftag{FOL}{\Sat{M}}{\pSat{v}}{\Gamma}$. In het bijzonder geldt
  $\iftag{FOL}{\Sat{M}}{\pSat{v}}{\lnot !B}$. Dus
  $\iftag{FOL}{\Sat/{M}}{\pSat/{v}}{!B}$, d.w.z.
  $\iftag{FOL}{\Struct{M}}{\pAssign{v}}$ vervult
  $\sFmla{\False}{!B}$.
\item De tak wordt uitgebreid door $\TRule{\False}{\lnot}$ toe te
  passen op $\sFmla{\False}{\lnot !B} \in \Gamma$: opgave.
\item De tak wordt uitgebreid door $\TRule{\True}{\land}$ toe te
  passen op $\sFmla{\True}{!B \land !C} \in \Gamma$. Dit levert twee
  nieuwe !!{signed formula}s op de tak op: $\sFmla{\True}{!B}$ en
  $\sFmla{\True}{!C}$. Stel dat
  $\iftag{FOL}{\Sat{M}}{\pSat{v}}{\Gamma}$, en in het bijzonder
  $\iftag{FOL}{\Sat{M}}{\pSat{v}}{!B \land !C}$. Dan
  $\iftag{FOL}{\Sat{M}}{\pSat{v}}{!B}$ en
  $\iftag{FOL}{\Sat{M}}{\pSat{v}}{!C}$. Dit betekent dat
  $\iftag{FOL}{\Struct{M}}{\pAssign{v}}$ zowel
  $\sFmla{\True}{!B}$ als $\sFmla{\True}{!C}$ vervult.
\item De tak wordt uitgebreid door $\TRule{\False}{\lor}$ toe te
  passen op $\sFmla{\False}{!B \lor !C} \in \Gamma$: opgave.
\item De tak wordt uitgebreid door $\TRule{\False}{\lif}$ toe te
  passen op $\sFmla{\False}{!B \lif !C} \in \Gamma$. Dit levert twee
  nieuwe !!{signed formula}s op de tak op: $\sFmla{\True}{!B}$ en
  $\sFmla{\False}{!C}$. Stel dat
  $\iftag{FOL}{\Sat{M}}{\pSat{v}}{\Gamma}$, en in het bijzonder
  $\iftag{FOL}{\Sat/{M}}{\pSat/{v}}{!B \lif !C}$. Dan
  $\iftag{FOL}{\Sat{M}}{\pSat{v}}{!B}$ en
  $\iftag{FOL}{\Sat/{M}}{\pSat/{v}}{!C}$. Dit betekent dat
  $\iftag{FOL}{\Struct{M}}{\pAssign{v}}$ zowel
  $\sFmla{\True}{!B}$ als $\sFmla{\False}{!C}$ vervult.

\iftag{FOL}{%
\item De tak wordt uitgebreid door $\TRule{\True}{\lforall}$ toe te
  passen op $\sFmla{\True}{\lforall[x][!B(x)]} \in \Gamma$. Dit levert
  op de tak een nieuwe !!{signed formula}~$\sFmla{\True}{!A(t)}$ op.
  Stel dat $\Sat{M}{\Gamma}$, en in het bijzonder
  $\Sat{M}{\lforall[x][!A(x)]}$. Volgens
  \olref[syn][sem]{prop:quant-terms} geldt $\Sat{M}{!A(t)}$.
  Bijgevolg vervult $\Struct{M}$ de formule
  $\sFmla{\True}{!A(t)}$.

\item De tak wordt uitgebreid door $\TRule{\False}{\lforall}$ toe te
  passen op $\sFmla{\False}{\lforall[x][!B(x)]} \in \Gamma$. Dit
  levert een nieuwe !!{signed formula}~$\sFmla{\False}{!A(a)}$ op,
  waarbij $a$ !!a{constant} is die niet in~$\Gamma$ voorkomt. Omdat
  $\Gamma$ vervulbaar is, bestaat er een $\Struct{M}$ zodat
  $\Sat{M}{\Gamma}$, en in het bijzonder
  $\Sat/{M}{\lforall[x][!B(x)]}$. We moeten aantonen dat
  $\Gamma \cup \{\sFmla{\False}{!A(a)}\}$ vervulbaar is. Daartoe
  definiëren we als volgt een geschikte~$\Struct{M'}$.

  Volgens \olref[syn][ass]{prop:sat-quant} geldt
  $\Sat/{M}{\lforall[x][!B(x)]}$ dan en slechts dan als voor een zekere
  $s$ geldt dat $\Sat/{M}{!B(x)}[s]$. Laat $\Struct{M'}$ nu gelijk zijn
  aan $\Struct{M}$, behalve dat $\Assign{a}{M'} = s(x)$. Volgens
  \olref[syn][ext]{cor:extensionality-sent} geldt voor elke
  $\sFmla{\True}{!C} \in \Gamma$ dat $\Sat{M'}{!C}$, en voor elke
  $\sFmla{\False}{!C} \in \Gamma$ dat $\Sat/{M'}{!C}$, omdat $a$ niet
  in~$\Gamma$ voorkomt.

  Volgens \olref[syn][ext]{prop:extensionality} geldt
  $\Sat/{M'}{!A(x)}[s]$. Volgens
  \olref[syn][ext]{prop:ext-formulas} geldt $\Sat/{M'}{!A(a)}[s]$.
  Omdat $!A(a)$ !!a{sentence} is, geldt volgens
  \olref[syn][ass]{prop:sentence-sat-true} dat
  $\Sat/{M'}{!A(a)}$, d.w.z. $\Struct{M'}$ vervult
  $\sFmla{\False}{!A(a)}$.
\item De tak wordt uitgebreid door $\TRule{\True}{\lexists}$ toe te
  passen op $\sFmla{\True}{\lexists[x][!B(x)]} \in \Gamma$: opgave.
\item De tak wordt uitgebreid door $\TRule{\False}{\lexists}$ toe te
  passen op $\sFmla{\False}{\lexists[x][!B(x)]} \in \Gamma$: opgave.}{}
\end{enumerate}
Nu beschouwen we de mogelijke inferenties die wel een gesplitste tak
opleveren.
\begin{enumerate}
\item De tak wordt uitgebreid door $\TRule{\False}{\land}$ toe te
  passen op $\sFmla{\False}{!B \land !C} \in \Gamma$. Dit levert twee
  takken op: een linkertak die via $\sFmla{\False}{!B}$ verdergaat en
  een rechtertak die via $\sFmla{\False}{!C}$ verdergaat. Stel dat
  $\iftag{FOL}{\Sat{M}}{\pSat{v}}{\Gamma}$, en in het bijzonder
  $\iftag{FOL}{\Sat/{M}}{\pSat/{v}}{!B \land !C}$. Dan
  $\iftag{FOL}{\Sat/{M}}{\pSat/{v}}{!B}$ of
  $\iftag{FOL}{\Sat/{M}}{\pSat/{v}}{!C}$. In het eerste geval vervult
  $\iftag{FOL}{\Struct{M}}{\pAssign{v}}$
  $\sFmla{\False}{!B}$, d.w.z.
  $\iftag{FOL}{\Struct{M}}{\pAssign{v}}$ vervult de formules op de
  linkertak. In het tweede geval vervult
  $\iftag{FOL}{\Struct{M}}{\pAssign{v}}$ $\sFmla{\False}{!C}$, d.w.z.
  $\iftag{FOL}{\Struct{M}}{\pAssign{v}}$ vervult de formules op de
  rechtertak.
\item De tak wordt uitgebreid door $\TRule{\True}{\lor}$ toe te
  passen op $\sFmla{\True}{!B \lor !C} \in \Gamma$: opgave.
\item De tak wordt uitgebreid door $\TRule{\True}{\lif}$ toe te
  passen op $\sFmla{\True}{!B \lif !C} \in \Gamma$: opgave.
\item De tak wordt met \Cut uitgebreid. Dit levert twee takken op:
  de ene bevat $\sFmla{\True}{!B}$ en de andere bevat
  $\sFmla{\False}{!B}$. Omdat $\iftag{FOL}{\Sat{M}}{\pSat{v}}{\Gamma}$
  en ofwel $\iftag{FOL}{\Sat{M}}{\pSat{v}}{!B}$, ofwel
  $\iftag{FOL}{\Sat/{M}}{\pSat/{v}}{!B}$,
  vervult $\iftag{FOL}{\Struct{M}}{\pAssign{v}}$ de linker- of de
  rechtertak.
\end{enumerate}
\end{proof}

\tagprob{FOL}
\begin{prob}
Voltooi het bewijs van \olref[fol][tab][sou]{thm:tableau-soundness}.
\end{prob}
\tagendprob

\tagprob{notFOL}
\begin{prob}
Voltooi het bewijs van \olref[pl][tab][sou]{thm:tableau-soundness}.
\end{prob}
\tagendprob

\begin{cor}
\ollabel{cor:weak-soundness}
Als $\Proves !A$, dan is $!A$ \iftag{FOL}{geldig}{een tautologie}.
\end{cor}

\begin{cor}
\ollabel{cor:entailment-soundness}
Als $\Gamma \Proves !A$, dan geldt $\Gamma \Entails !A$.
\end{cor}

\begin{proof}
  Als $\Gamma \Proves !A$, dan heeft voor zekere $!B_1$, \dots,
  $!B_n \in \Gamma$ de verzameling
  $\{\sFmla{\False}{!A}, \sFmla{\True}{!B_1}, \dots,
  \sFmla{\True}{!B_n}\}$ een gesloten !!{tableau}. Volgens
  \olref{thm:tableau-soundness} maakt elke
  \iftag{FOL}{!!{structure}}{!!{valuation}}~$\iftag{FOL}{\Struct{M}}{\pAssign{v}}$
  een zekere $!B_i$ onwaar of $!A$ waar. Als dus
  $\iftag{FOL}{\Sat{M}}{\pSat{v}}{\Gamma}$, dan ook
  $\iftag{FOL}{\Sat{M}}{\pSat{v}}{!A}$.
\end{proof}

\begin{cor}
\ollabel{cor:consistency-soundness}
Als $\Gamma$ vervulbaar is, dan is zij consistent.
\end{cor}

\begin{proof}
We bewijzen de contrapositie. Stel dat $\Gamma$ niet consistent is.
Dan zijn er $!B_1$, \dots, $!B_n \in \Gamma$ en een gesloten
!!{tableau} voor $\{\sFmla{\True}{!B_1}, \dots, \sFmla{\True}{!B_n}\}$.
Volgens \olref{thm:tableau-soundness} is er geen
$\iftag{FOL}{\Struct{M}}{\pAssign{v}}$ waarvoor
$\iftag{FOL}{\Sat{M}}{\pSat{v}}{!B_i}$ voor alle $i=1$, \dots,~$n$.
Maar dan is $\Gamma$ niet vervulbaar.
\end{proof}

\end{document}
